% Section 3: Methods (Compressed Version)
% Focus: Reproducible pipeline, experimental setup, training details

\section{Experimental Methodology}
\label{sec:methods}

We conduct 153 systematically controlled experiments to understand ternary quantization for CNNs. All experiments follow a reproducible pipeline from code execution to paper generation, enabling complete verification of our claims.

\subsection{Reproducible Research Pipeline}
\label{sec:methods:pipeline}

Our end-to-end pipeline eliminates manual transcription through automated artifact generation:

\begin{enumerate}
    \item \textbf{Training:} 153 experiments produce \texttt{results.json} (metrics), \texttt{config.json} (hyperparameters), and \texttt{best\_model.pth} (checkpoint)
    \item \textbf{Aggregation:} \texttt{aggregate\_results.py} loads all results into \texttt{aggregated.csv}
    \item \textbf{Generation:} Scripts produce LaTeX tables and PDF figures directly from CSV
    \item \textbf{Paper:} LaTeX commands directly reference generated tables (\verb|\input{tables/*.tex}|) and figures (\verb|\includegraphics{figures/*.pdf}|)
\end{enumerate}

Every number in this paper traces directly to experimental data with documented parameters (see Appendix~\ref{app:reproducibility} for complete pipeline).

\subsection{Datasets and Architectures}

\textbf{Datasets:} CIFAR-10 (10 classes, 32×32), CIFAR-100 (100 classes, 32×32), Tiny-ImageNet (200 classes, 64×64)~\cite{krizhevsky2009cifar, le2015tinyimagenet}. Standard splits used without modifications.

\textbf{Architectures:} ResNet-18 (11.17M params) and ResNet-50 (23.52M params)~\cite{he2016resnet} with \textbf{CIFAR-adapted stems}: 3×3 conv (stride 1), no maxpool. This preserves spatial resolution (32×32 → 32×32) and matches standard CIFAR practice~\cite{pytorch_cifar_kuangliu, pytorch_cifar100_weiaicunzai}. Section~\ref{sec:architecture_choice} shows this recovers +17\% accuracy on CIFAR-100 vs standard ImageNet stems.

\subsection{Training Configurations}

Table~\ref{tab:training_config} summarizes our training setup. All experiments use identical base configuration with variations only in quantization and knowledge distillation.

\begin{table}[h]
\centering
\caption{Training configuration (applies to FP32, BitNet, and TTQ experiments).}
\label{tab:training_config}
\begin{tabular}{ll}
\toprule
\textbf{Parameter} & \textbf{Value} \\
\midrule
\textbf{Optimizer} & SGD (momentum 0.9, weight decay 5e-4) \\
\textbf{Learning rate} & 0.1, cosine annealing to 0, warmup 5 epochs \\
\textbf{Batch size} & 128 (64 for R50 Tiny-ImageNet) \\
\textbf{Epochs} & 200 (baseline), 300 (recipe experiments) \\
\textbf{Data augmentation} & Random crop + horizontal flip (basic) \\
\textbf{Advanced augmentation} & Mixup ($\alpha$=0.2), label smoothing ($\epsilon$=0.1) \\
                               & (CIFAR-10, Tiny-ImageNet only; not CIFAR-100) \\
\textbf{BitNet quantization} & Weights: $\{-1, 0, +1\}$ via \texttt{sign()} \\
                             & Activations: FP16 (no quantization) \\
\textbf{Knowledge Distillation} & Temperature T=4.0, $\alpha$=0.9 \\
                                & Teacher: FP32 model (seed 42) \\
\midrule
\textbf{Hardware} & NVIDIA A100/V100 GPUs \\
\textbf{Total compute} & 918 GPU-hours (153 experiments) \\
\textbf{Seeds} & 42, 123, 456 (3 per configuration) \\
\bottomrule
\end{tabular}
\end{table}

\subsection{Trained Ternary Quantization (TTQ)}

To address Reviewer 2's requirement for comparison with state-of-the-art ternary quantization, we include Trained Ternary Quantization (TTQ)~\cite{zhu2017ttq} experiments. Zhu et al. demonstrate that learning per-layer positive and negative scaling factors ($w_p$, $w_n$) alongside a quantization threshold ($\delta$) achieves near-zero accuracy gaps on CIFAR-10, substantially improving upon fixed ternary schemes.

TTQ quantizes weights to $\{-w_n, 0, +w_p\}$ where the scales and threshold are learned via gradient descent. During forward propagation, weights $w$ are quantized based on the learned threshold: $w > \delta$ maps to $+w_p$, $w < -\delta$ maps to $-w_n$, and $|w| \leq \delta$ maps to $0$. The scales and threshold are initialized from weight statistics and refined during training using custom autograd functions to ensure gradient flow. Unlike BitNet's fixed $\{-1, 0, +1\}$ scheme, TTQ adds 2 FP32 parameters per layer but enables asymmetric scales that adapt to layer-specific weight distributions.

We train TTQ models using identical hyperparameters to our Phase 1 FP32 baselines (300 epochs, SGD with lr=0.1, cosine schedule, mixup and label smoothing on CIFAR-10/Tiny-ImageNet). This matched-condition setup enables fair comparison: any accuracy differences reflect quantization scheme effectiveness rather than training recipe differences. TTQ activations remain FP32 (matching our BitNet implementation), isolating the comparison to ternary weight quantization strategies. All 18 TTQ experiments converged stably with no observed training instabilities or parameter divergence.

\subsection{Experiment Matrix}

Our 153 experiments comprise 135 original experiments plus 18 TTQ comparisons (2 architectures × 3 datasets × 3 seeds). The original 135 decompose into 5 baseline types:

\begin{table}[h]
\centering
\caption{Experimental design: 135 baseline experiments + 18 TTQ comparisons = 153 total.}
\label{tab:experiment_matrix}
\begin{tabular}{lp{10cm}}
\toprule
\textbf{Baseline} & \textbf{Purpose} \\
\midrule
\textbf{FP32} (18 exp) & Standard full-precision baseline \\
\textbf{FP32+KD} (9 exp) & Isolate KD benefit independent of quantization (critical control) \\
\textbf{BitNet} (18 exp) & Pure ternary quantization (no KD, no mixed-precision) \\
\textbf{BitNet+KD} (9 exp) & Ternary quantization with knowledge distillation only \\
\textbf{BitNet+KD+keep\_conv1} (18 exp) & Full recipe: KD + FP32 conv1 + ternary elsewhere \\
\midrule
\textbf{Layer ablations} (63 exp) & Keep specific layers in FP32 (conv1, layer1, layer4, fc) \\
                                   & to measure per-layer sensitivity \\
\bottomrule
\midrule
\textbf{TTQ} (18 exp) & Trained Ternary Quantization comparison with learned per-layer scales \\
\bottomrule
\multicolumn{2}{l}{\footnotesize Phase 1 baseline: 135 = 18+9+18+9+18+63 across (R18, R50) × (CIFAR-10, CIFAR-100, Tiny-ImageNet) × (s42, s123, s456)} \\
\multicolumn{2}{l}{\footnotesize Phase 6 TTQ: 18 = (R18, R50) × (CIFAR-10, CIFAR-100, Tiny-ImageNet) × (s42, s123, s456)}
\end{tabular}
\end{table}

Our experimental design ensures fair comparison through four principles. First, we include FP32+KD baselines to isolate quantization effects from training improvements, addressing a common confound in quantization literature. Second, all models use identical training recipes (optimizer, learning rate, schedule, augmentation) with only quantization method and KD as experimental variables. Third, we maintain statistical rigor through 3 seeds per configuration with mean±std reporting and paired t-tests for key comparisons. Fourth, we deliberately use basic augmentation (random crop + flip) rather than heavy augmentation—prior exploratory work suggested that techniques like cutout, RandAugment, and mixup widened FP32-ternary gaps by 1.5-2.5×, making them unsuitable for fair baseline comparison.

\subsection{Deterministic Training}

All experiments are fully reproducible with fixed seeds:
\begin{itemize}
    \item \textbf{PyTorch:} \texttt{torch.manual\_seed(seed)}, \texttt{torch.backends.cudnn.deterministic = True}
    \item \textbf{NumPy:} \texttt{np.random.seed(seed)}
    \item \textbf{DataLoader:} Fixed \texttt{worker\_init\_fn} for deterministic augmentation
\end{itemize}

Same seed + same code = bit-exact identical checkpoints (verified via MD5 hash).

\subsection{Statistical Analysis}

For each configuration, we report:
\begin{itemize}
    \item \textbf{Mean ± standard deviation} across 3 seeds
    \item \textbf{Accuracy gap:} FP32 - BitNet (baseline gap), FP32+KD - BitNet+Recipe (final gap)
    \item \textbf{Gap recovery:} Percentage of baseline gap closed by recipe
    \item \textbf{Statistical tests:} Two-sample t-tests (p-values) and Cohen's d (effect sizes)
\end{itemize}

All metrics computed programmatically from \texttt{aggregated.csv} using \texttt{generate\_tables.py} and \texttt{generate\_figures.py}.
